\documentclass{ximera}

% MATH -----------------------------------------------------------
\newcommand{\nats}{\mathbb N}
\newcommand{\ints}{\mathbb Z}
\newcommand{\rats}{\mathbb Q}
\newcommand{\reals}{\mathbb R}
\newcommand{\complex}{\mathbb C}
\newcommand{\powerset}{\mathscr P}

%-------------------------------------------------------------

\pgfplotsset{compat=1.18}
\begin{document}
\begin{abstract}
 \end{abstract}

    \title{Exemplar Examples 19}
    
    \maketitle

Suppose an RLC circuit, with a resistor of resistance $R=100$ ohms, an inductor of inductance $L=0.5$ henries, and a capacitor with a capacitance of \\$C=2\times 10^{-4}$ farads, has (at time $t=0$) an initial current of $25$ amps that is initially decreasing at $1$ amp per second, and is attached to a 9-volt battery.

\begin{enumerate}

\item Let's derive an IVP satisfied by the current $I(t)$ at $t\geq 0$ seconds.

\vfill


% Note: V(t)=9 so V'(t)=0 and

% 0.5 I''+00I'+5000I=0, I(0)=25, I'(0)=-1.

% I''+200I'+10,000I=0, I(0)=25, I'(0)=-1.


\item Let's solve this IVP.

\vfill
\vfill


% r^2+200r+10000 =(r+100)^2 ---> I=(c_1+c_2 t)e^{-100t}


% I=(25+c_2 t)e^{-100t}

% I'=(c_2-100(25+c_2 t))e^{-100t} ---> -1=c_2-2500 ---> c_2=2499

% I=(25+2499t)e^(-100t)


\newpage

\item What happens to the current in the long run?  And what happens to the (positive) charge on one of plates of the capacitor in the long run?


% It dissipates towards 0.  It pushes charge on to the capacitor:  9 Volts = Q/C  so Q = 9*2\times 10^{-4} = 1.8\times 10^{-3} coulombs



\end{enumerate}

\end{document}